\section{Conclusion and Future Work}
\label{sec:conclusion}

In this work, we introduced \textbf{Gaussian Process Dynamic Routing (GP-DR)}, a mathematically rigorous, non-parametric Bayesian dynamic routing framework designed to resolve the dual challenges of data scarcity and stream-level heterogeneity in modular deep learning. 

By modeling parameter routing functions under a Gaussian Process prior, we bypass the standard parametric training loop entirely, effectively neutralizing the Overfitting-Optimizer Paradox. Under extreme calibration constraints ($N=64$ samples), GP-DR achieves a $72.40\%$ test Joint Mean accuracy on the Isolating Coordinate Sandbox with zero training loops. While the predictive posterior variance theoretically provides a mathematically bounded metric of epistemic uncertainty to flag out-of-distribution (OOD) inputs, our analysis reveals a critical unit-sphere variance collapse under realistic noise and that simpler distance heuristics outperform it under representational overlap. Furthermore, we demonstrated that integrating stream-level \textbf{Micro-Batch Homogenization (MBH)} completely resolves vectorization collapse under heterogeneous streams, enabling a $70.20\%$ recovery under mixed-task batching.

\textbf{Future Directions:}
There are several promising avenues for future research. First, we plan to explore and deploy alternative positive-definite kernels—specifically Cosine/Inner-Product and Mahalanobis kernels (formulated in Section 3.5)—as a primary research direction, bypassing the Euclidean RBF kernel's lengthscale boundaries and natively supporting anisotropic densities under representation-space projection. Along this direction, we will study adaptive kernel learning on representation manifolds, allowing lengthscale and noise parameters to adapt dynamically across layers. Second, we aim to scale GP-DR to larger, mid-sized and massive architectures (such as RoBERTa-Base or LLaMA-3B/8B generative models) where representation dimensions are substantially higher (e.g., $D \in \{768, 1024, 4096\}$). Under massive representational manifolds, we expect the concentration of measure and high-dimensional representational sparsity to work in our favor. In extremely high-dimensional vector spaces, the number of orthogonal directions scales linearly with $D$, which mathematically guarantees that task-specialized representations reside in naturally near-orthogonal subspaces. Consequently, the subspace coordinate projection becomes structurally more robust and less susceptible to representational blur in massive backbones. Evaluating GP-DR on these massive architectures will allow us to empirically verify if the high-dimensional representational sparsity naturally mitigates GPR task-conflict spatial blindspots without requiring extreme subspace coordinate projections, facilitating a more direct, non-parametric Bayesian integration over massive pre-trained model pools. For settings with severe manifold overlap, we propose localized GP routing (e.g., neighborhood-based sparse GPR) or our formulated Mahalanobis and Cosine kernels (Section 3.5) to warp the metric space, natively resolving any potential task-conflict blindspots. Finally, we intend to explore non-Gaussian likelihood models and variational inference to provide a true categorical Bayesian framework that natively respects probability simplex constraints while maintaining real-time online inference capabilities.
